\begin{enumerate}[label=\thesubsection.\arabic*,ref=\thesubsection.\theenumi]
    \item A spring-mass-damper system with a mass of $1 \, \mathrm{kg}$ is found to have a damping ratio of 0.2 and a natural frequency of $5 \, \mathrm{rad/s}$. The damping of the system is given by 
	    \hfill (AE 2007)
    \begin{multicols}{4}
    \begin{enumerate}
        \item $2 \, \mathrm{Ns/m}$ 
        \item $2 \, \mathrm{N/s}$ 
        \item $0.2 \, \mathrm{kg/s}$ 
        \item $0.2 \, \mathrm{N/m}$
    \end{enumerate}
    \end{multicols}
    \item A 1 kg mass attached to a spring elongates it by 16 mm. The mass is then pulled from its equilibrium position by 10 mm and released from rest. Assuming the acceleration due to gravity of $9.81 m/s^2$, the response of the mass in mm is given by: 
    \hfill (AE 2007)
\begin{enumerate}
    \begin{multicols}{4}
        \item $x = \cos 4.76 \, t$
        \item $x = \sin 4.76 \, t$
        \item $x = \sin 16 \, t$
        \item $x = 10 \cos 16 \, t$ 
	\end{multicols}
    \end{enumerate}
    \item A spring-mass-damper system is excited by a force $F_0 \sin \omega t$. The amplitude at resonance is measured to be $1 \, \mathrm{cm}$. At half the resonance frequency, the amplitude is 0.5 cm. The damping ratio of the system is 
    \hfill (AE 2007)
    \begin{multicols}{4}
\begin{enumerate}
        \item $0.1026$ \quad
        \item $0.3242$ \quad
        \item $0.7211$ \quad
        \item $0.1936$
    \end{enumerate}
    \end{multicols}

    \item $\mathrm{\lim_{x \to 0} \frac{\sin x}{e^4 x} =}$
    \hfill (AE 2007)
    \begin{multicols}{4}
\begin{enumerate}
        \item 10
        \item 0
        \item 1
        \item $\infty$
    \end{enumerate}
    \end{multicols}
    \item Let a dynamical system be described by the differential equation $2 \frac{dx}{dt} + \cos x = 0$. Which of the following differential equations describes this system in a close approximation sense for small perturbation about $x = \pi/4$?
    \hfill (AE 2007)
    \begin{multicols}{2}
\begin{enumerate}
        \item $2 \frac{dx}{dt} + \sin x = 0$
        \item $2 \frac{dx}{dt} - \frac{1}{\sqrt{2}} x = 0$
        \item $\frac{dx}{dt} + \cos x = 0$
        \item $\frac{dx}{dt} + x = 0$
    \end{enumerate}
    \end{multicols}
\item  Consider the spring mass system shown in the figure below. This system has two degrees of freedom representing the motions of the two masses.
\hfill (AE 2007)

\begin{figure}[H]
    \centering
    \input{figs/q74-75.tex}
    \caption{}
    \label{fig:q74 75}
\end{figure}

\begin{enumerate}
    \item The system shows the following type of coordinate coupling
    \begin{enumerate}
        \item Static coupling
        \item Dynamic coupling
        \item Static and dynamic coupling
        \item No coupling
    \end{enumerate}

    \item The two natural frequencies of the system are given as
    \begin{multicols}{4}
    \begin{enumerate}
        \item $\sqrt{\frac{4 \pm \sqrt{5} k}{3 m}}$
        \item $\sqrt{\frac{4 \pm \sqrt{3} k}{3 m}}$
        \item $\sqrt{\frac{4 \pm \sqrt{7} k}{3 m}}$
        \item $\sqrt{\frac{4 \pm \sqrt{11} k}{3 m}}$
    \end{enumerate}
    \end{multicols}
    
\end{enumerate}

\item Let $$F(s) = \frac{(s + 10)}{(s +2)(s + 20)}$$

\hfill (AE 2007)
\begin{enumerate}
    \item The partial fraction expansion of $F(s)$ is
    \begin{multicols}{4}
    \begin{enumerate}
        \item $\frac{1}{s + 2} + \frac{1}{s + 20}$
        \item $\frac{5}{s + 2} + \frac{2}{s + 20}$
        \item $\frac{2}{s + 2} + \frac{20}{s + 20}$
        \item $\frac{4/9}{s + 2} + \frac{5/9}{s + 20}$
    \end{enumerate}
    \end{multicols}
    \item The inverse Laplace transform of $F(s)$ is
    \begin{multicols}{2}
    \begin{enumerate}
        \item $2e^{-34} + 20e^{-34}$
        \item $\frac{4}{9}e^{-34} + \frac{5}{9}e^{-34}$
        \item $5e^{-34} + 2e^{-34}$
        \item $\frac{9}{4}e^{-34} + \frac{9}{5}e^{-34}$
    \end{enumerate}
    \end{multicols}
    
\end{enumerate}
\item 
\begin{align*}
\frac{\omega}{(s + a)^2 + \omega^2} &\quad 
\end{align*}
is the Laplace Transform of
\hfill(AG 2007)
\begin{multicols}{2}
\begin{enumerate}
  
  \item $\exp(-at) \sin \omega t$
  
  \item $\sin \omega t$
  
   \item $\exp(-at) \sinh \omega t$
  
  \item $\sinh \omega t$
\end{enumerate}
\end{multicols}

\end{enumerate}

